\chapter{Why the interconnect matters}
\label{ch:role}

\section{The two themes of this book}

Large AI systems fail or succeed as much on wiring as on FLOPs. The short-reach
optics that connect accelerators, switches, and memory sit on the critical path
for latency, power, and uptime, and two parts of that layer decide whether the
system scales: the \term{lasers} that generate the light, and the discipline of
\term{IM/DD validation}\sidenote{IM/DD = intensity modulation with direct
detection: the short-reach datacenter transceiver style. See \Cref{ch:imdd}.}
that proves those links work from the first bench eye through production and
fleet life. This book keeps returning to those two themes. Modulation formats,
WDM, reliability, and network architecture matter because they serve them.

\keyidea{Two things dominate short-reach optics at scale: \emph{lasers} and
\emph{IM/DD validation}. Everything else orbits those two.}

\section{The context: purpose-built inference silicon}

The pressure on the interconnect is not abstract. It comes from a shift in how
large AI systems are built: vertically integrated, purpose-built silicon
deployed at gigawatt scale, with networking named as a first-order design axis
beside compute and memory. A representative example is \term{Jalapeño}, a purpose-built LLM \emph{inference}
accelerator Broadcom and a hyperscaler partner announced in 2026 as a blank-slate
design, the first chip in a multi-generation compute
platform.\sidenote{Source: \citehref{jalapeno2026}{Broadcom and partner, ``LLM-optimized
inference chip,'' 24 June 2026}.} Its public
features are typical of the class:

\begin{itemize}
\item \textbf{Partners and integration.} Broadcom provides silicon implementation
      and networking (including \term{Tomahawk} switch silicon); Celestica
      provides board, rack, and system integration.
\item \textbf{Cadence.} A roughly nine-month design-to-tape-out cycle, with the
      designers' own models used to accelerate parts of the flow.
\item \textbf{Scale.} Deployment at gigawatt scale, over multiple generations,
      beginning in late 2026.
\item \textbf{Design thesis.} ``Reduce data movement and balance compute, memory,
      and \emph{networking} resources'' to reach realized utilization close to
      theoretical peak.
\end{itemize}

That last point is the hook for this book. Purpose-built inference silicon does
not treat networking as a commodity NIC bolted on after the die is done; it
budgets interconnect power and latency alongside FLOPs and HBM. Partners and
cadence matter because they set who owns the switch ASIC and how fast generations
turn: Broadcom for silicon and networking (including Tomahawk), Celestica for
board and rack, a roughly nine-month design-to-tape-out cycle, and gigawatt-scale
deployment planned across multiple generations from late 2026. Once networking
sits on that line, the optical interconnect is how the system scales past a
single package to a gigawatt cluster, and laser quality plus IM/DD validation
become infrastructure problems rather than module afterthoughts.

\section{Why inference makes the interconnect matter}

Inference is not training. Once a model is trained, serving it is dominated by
two phases with very different bottlenecks (developed fully in
\Cref{ch:networking}):

\begin{description}
\item[Prefill] processing the prompt: highly parallel and compute-bound, much
      like training.
\item[Decode] generating tokens one at a time: autoregressive and
      \emph{memory-bandwidth-bound}, because every token streams the model
      weights through the compute units again.
\end{description}

Frontier models do not fit on one chip. They are sharded across many
accelerators, so every generated token triggers \term{collective communication}
across the fabric: all-reduce for tensor parallelism, all-to-all for
mixture-of-experts routing, point-to-point for pipeline stages. The interconnect
therefore sits on the \emph{latency critical path} of inference, not merely the
plumbing between training runs.

\keyidea{Reliable, low-power IM/DD links directly set how large and how
dependable an inference fabric can be. That is why the optical layer has become
central to AI system design.}

\section{How to read this book}

The chapters build from physics to fleet scale:

\begin{enumerate}
\item \Cref{ch:firstprinciples,ch:imdd}: energy, IM/DD vocabulary, modulators, FEC, equalization.
\item \Cref{ch:models}: quantitative noise, RIN, sensitivity (use with \Cref{sec:link-budget}).
\item \Cref{ch:lasers}: light sources (DFB/EML, LIV/SMSR/RIN, aging, ELSFP/CW-WDM).
\item \Cref{ch:wdm}: wavelength locking, thermal crosstalk, CW-WDM, on-chip MUX.
\item \Cref{ch:validation,ch:reliability}: measurement ladder, link budgets, qual, packaging.
\item \Cref{ch:networking}: scale-up/out, pluggables, CPO/XPO, inference collectives.
\end{enumerate}

To use the book as a design drill, pick one link style (retimed 800G DR, LPO, or
CPO WDM) and trace it end to end through \Cref{sec:txrx-chain,sec:pluggables,sec:cpo-status}.
